\section{Introduction}

Ce document fait le point sur le lancement du projet : ce que l'équipe a compris du besoin du Cerema, ce qu'elle va réaliser et comment elle s'organise. Il s'adresse au client et au tuteur école, prépare la revue de projet et sert à affiner le besoin avec le client. Il s'appuie sur la fiche du sujet, les travaux du projet 2025, la réunion de lancement et les échanges avec le Cerema (section~\ref{sec:bibliographie}). Le tableau~\ref{tab:terminologie} définit les termes techniques utilisés.

\begin{table}[H]
\centering
\small
\caption{Terminologie}
\label{tab:terminologie}
\begin{tabularx}{\textwidth}{@{}>{\raggedright\arraybackslash}p{3.8cm}L@{}}
\toprule
\textbf{Terme} & \textbf{Définition} \\
\midrule
USV (\emph{Unmanned Surface Vehicle}) & Drone de surface : bateau sans équipage, télépiloté ou autonome. \\
PCE & Projet Commande Entreprise, réalisé par des élèves pour un partenaire extérieur. \\
Autopilote (Pixhawk, ArduPilot) & Carte et logiciel qui commandent les moteurs et gèrent la navigation. \\
Station sol (QGroundControl) & Ordinateur qui permet de piloter le bateau, de programmer et de suivre les missions. \\
Ordinateur compagnon & Ordinateur embarqué (ici un Raspberry Pi 4) qui analyse les images et dialogue avec l'autopilote. \\
MAVLink & Langage de communication entre l'autopilote, l'ordinateur compagnon et la station sol. \\
Point de passage (\emph{waypoint}) & Position que le bateau doit rejoindre lors d'une mission programmée. \\
ESC & Variateur qui commande la vitesse d'un moteur. \\
YOLO & Famille de modèles d'IA qui détectent des objets dans une image. \\
Métadonnées de détection & Résultats du modèle (cadre, classe, score de confiance) transmis avec l'image, exploitables par l'ordinateur compagnon. \\
OCR & Reconnaissance de caractères, ici pour lire l'immatriculation d'un navire. \\
AIS & Système par lequel un navire émet en continu son identité, sa position et sa vitesse. \\
Gouverne & Surface orientable qui dirige le bateau, ici commandée par un moteur. \\
\bottomrule
\end{tabularx}
\end{table}
